\subsection{\texorpdfstring{\(\mathfrak{S}\)-收敛}{S-收敛}}

定义 2 。设 X 是一个集，Y 是一致空间，S 是 X 的子集组成的一个集。在 S 的集上一致收敛的一致结构，或简称 S-收敛的一致结构，是 F(X; Y) 上使限制映射 \(u \to u|A\)（F(X, Y) 到一致空间 \(\mathcal{F}_u(A; Y)\) 中的，其中 A 取遍 S）一致连续的最粗一致结构。把 F(X; Y) 赋以 S-收敛的一致结构所得的一致空间记作 \(\mathcal{F}_{\mathfrak{S}}(X; Y)\)。

\(\mathfrak{S}\)-收敛的一致结构所诱导的拓扑称为 \(\mathfrak{S}\)-收敛的拓扑；它是使所有映射 \(u \to u|A\)（\(\mathcal{F}(X; Y)\) 到 \(\mathcal{F}_u(A; Y) (A \in \mathfrak{S})\) 中的）都连续的最粗拓扑（第 II 章，§ 2，第 3 号，命题 4 的推论）。

滤子 \(\Phi\)（\(\mathcal{F}(\mathbf{X};\mathbf{Y})\) 上的）收敛于 \(u_{0}\)（关于 \(\mathfrak{S}\)-收敛的拓扑），当且仅当 \(u|\mathrm{A}\) 一致收敛于 \(u_{0}|\mathrm{A}\)（关于 \(\Phi\)，对所有 \(\mathrm{A}\in \mathfrak{S}\)）（第 I 章，§ 7，第 6 号，命题 10），因此称 \(\Phi\) 一致收敛于 \(u_{0}\)（在 \(\mathfrak{S}\) 的集上）。

滤子基 \(\Phi\)（\(\mathcal{F}_{\mathfrak{S}}(\mathbf{X};\mathbf{Y})\) 上的）是柯西滤子基，当且仅当对每个 \(\mathbf{A}\in\mathfrak{S}\)，\(\Phi\) 在映射 \(u\to u|\mathbf{A}\) 下的像是 \(\mathcal{F}_{u}(\mathbf{A};\mathbf{Y})\) 上的柯西滤子基（第 II 章，§ 3，第 1 号，命题 4）。

设 \(f\) 是拓扑（相应地，一致）空间 \(Z\) 到 \(\mathcal{F}_{\mathfrak{S}}(X; Y)\) 中的映射。则 \(f\) 连续（相应地，一致连续），当且仅当对每个 \(A \in \mathfrak{S}\)，映射 \(z \to f(z)|A\)（\(Z\) 到 \(\mathcal{F}_u(A; Y)\) 中的）连续（相应地，一致连续）（第 I 章，§ 2，第 3 号，命题 4；第 II 章，§ 2，第 3 号，命题 4）。

最后，设 \(\mathbf{M}\) 是 \(\mathcal{F}_{\mathfrak{S}}(\mathbf{X};\mathbf{Y})\) 的子集；则 \(\mathbf{M}\) 是准紧的，当且仅当对每个 \(\mathbf{A} \in \mathfrak{S}\)，限制 \(u|\mathbf{A}\)（\(u \in \mathbf{M}\)）组成的集是 \(\mathcal{F}_u(\mathbf{A};\mathbf{Y})\) 的准紧子集（第 II 章，§ 4，第 2 号，命题 3）。

注记。1) 初始一致结构的围域的一般定义（第 II 章，§ 2，第 3 号，命题 4）表明，\(\mathcal{F}_{\mathfrak{S}}(\mathrm{X};\mathrm{Y})\) 的围域基本系可如下得到：对每个 \(\mathbf{A}\in \mathfrak{S}\) 与每个围域 \(\mathbf{V}\)（\(\mathfrak{B}\) 的，此基本系属于 \(\mathbf{Y}\)），用 \(W(\mathbf{A},\mathbf{V})\) 表示映射的对 \((u,v)\)（\(\mathbf{X}\) 到 \(\mathbf{Y}\) 中的）中使 \((u(x),v(x))\in \mathbf{V}\)（对每个 \(x\in \mathbf{A}\)）者组成的集；当 \(\mathbf{A}\) 取遍 \(\mathfrak{S}\) 且 \(\mathbf{V}\) 取遍 \(\mathfrak{B}\) 时，诸集 \(W(\mathbf{A},\mathbf{V})\) 的有限交构成 \(\mathcal{F}_{\mathfrak{S}}(\mathrm{X};\mathrm{Y})\) 的围域基本系。

这描述立即表明：若 \(\mathfrak{S},\mathfrak{S}'\) 是 \(\mathbf{X}\) 的子集的两个集，使 \(\mathfrak{S} \subset \mathfrak{S}'\)，则 \(\mathfrak{S}'\)-收敛的一致结构比 \(\mathfrak{S}\)-收敛的一致结构更细。

\begin{enumerate}
\def\labelenumi{\arabic{enumi})}
\setcounter{enumi}{1}
\tightlist
\item
  然而，把 \(\mathfrak{S}\) 换成由含于 \(\mathfrak{S}\) 的集的有限并之中的 X 的所有子集组成的集 \(\mathfrak{S}'\)，\(\mathfrak{S}\)-收敛的一致结构不变。因此在研究 \(\mathfrak{S}\)-收敛时，我们总可以限于集 \(\mathfrak{S}\) 满足下列两个条件的情形：
\end{enumerate}

\((\mathbf{F}_{\mathbf{I}}^{\prime})\) \(\mathfrak{S}\) 中的集的每个子集都属于 \(\mathfrak{S}\)。

\((F_{\mathrm{II}}^{\prime})\) \(\mathfrak{S}\) 中的集的每个有限并都属于 \(\mathfrak{S}\)。

若 \((\mathbf{F}_{\mathrm{II}}^{\prime})\) 满足，则 \(\mathcal{F}_{\mathfrak{S}}(\mathbf{X};\mathbf{Y})\) 的一个围域基本系可这样得到：取所有集 \(\boldsymbol{W}(\mathrm{A},\mathrm{V})\)（其中 A 取遍 S，V 取遍 Y 的一个围域基本系）。

\begin{enumerate}
\def\labelenumi{\arabic{enumi})}
\setcounter{enumi}{2}
\tightlist
\item
  \(\mathfrak{S}\)-收敛的一致结构是这乘积空间的一致结构在映射 \(u \to (u|\mathbf{A})_{\mathbf{A} \in \mathfrak{S}}\)（\(\mathcal{F}(\mathbf{X}; \mathbf{Y})\) 到 \(\prod_{\mathbf{A} \in \mathfrak{S}} \mathcal{F}_u(\mathbf{A}; \mathbf{Y})\) 中的）下的逆像（第 II 章，§ 2，第 6 号，命题 8）。若 \(\mathfrak{S}\) 是 \(\mathbf{X}\) 的覆盖，则此映射是单射，因而 \(\mathcal{F}_{\mathfrak{S}}(\mathbf{X}; \mathbf{Y})\) 同构于 \(\prod_{\mathbf{A} \in \mathfrak{S}} \mathcal{F}_u(\mathbf{A}; \mathbf{Y})\) 的、作为此映射之像的一致子空间。
\end{enumerate}

命题 1 。若 Y 是豪斯多夫的且 S 是 X 的覆盖，则空间 \(\mathcal{F}_{\mathfrak{S}}(X;Y)\) 是豪斯多夫的。

设 \(u, v\) 是 \(\mathcal{F}_{\mathfrak{S}}(\mathrm{X}; \mathrm{Y})\) 的两个元素，使 \((u, v) \in W(\mathrm{A}, \mathrm{V})\)（对每个围域 \(\mathrm{V}\)（\(\mathrm{Y}\) 的）与每个 \(\mathrm{A} \in \mathfrak{S}\)）。由于 \(\mathrm{Y}\) 是豪斯多夫的，推出 \(u\) 与 \(v\) 在每个集 \(\mathrm{A} \in \mathfrak{S}\) 上一致；又因 \(\mathfrak{S}\) 覆盖 \(\mathrm{X}\)，必有 \(u = v\)。

注记。4) 设 H 是 \(\mathcal{F}(\mathbf{X};\mathbf{Y})\) 的子集。滥用术语，把 \(\mathfrak{S}\)-收敛的一致结构（相应地，拓扑）（\(\mathcal{F}(\mathbf{X};\mathbf{Y})\) 上的）在 H 上诱导的一致结构（相应地，拓扑）称为集 H 上 \(\mathfrak{S}\)-收敛的一致结构（相应地，拓扑）。

\begin{enumerate}
\def\labelenumi{\arabic{enumi})}
\setcounter{enumi}{4}
\tightlist
\item
  设 L 是由滤子 \(\mathfrak{G}\) 滤化的集，且 \(\lambda \to u_{\lambda}\) 是 L 到 \(\mathcal{F}_{\mathfrak{S}}(\mathbf{X};\mathbf{Y})\) 中的、有极限 \(v\)（关于 \(\mathfrak{G}\)）的映射；这时我们说，关于滤子 \(\mathfrak{G}\)，X 到 Y 中的映射 \(u_{\lambda}\) 一致收敛于 \(v\){[}或族 \((u_{\lambda})\) 一致收敛于 \(v\){]}（在 \(\mathfrak{S}\) 的每个集上）。若 \(\mathbf{L} = \mathbf{N}\) 且 \(\mathfrak{G}\) 是弗雷歇滤子，则在这陈述中省略对 \(\mathfrak{G}\) 的提及。
\end{enumerate}

更特殊地，设 Y 上定义了一个交换且结合的、以加法记号书写的合成律。若 \((u_{n})\) 是 X 到 Y 中的映射的任一序列，用 \(v_{n}\) 表示由下式定义的映射

\[
v _ {n} (x) = \sum_ {k = 0} ^ {n} u _ {k} (x) \quad (n \in \mathbf {N});
\]

我们说以 \(u_n\) 为通项的级数在 \(\mathfrak{S}\) 的每个集上一致收敛，如果序列 \((v_n)\) 在 \(\mathfrak{S}\) 的每个集上一致收敛。类似地，映射的一致可和族 \((u_\lambda)_{\lambda \in \mathbf{L}}\)（\(\mathbf{X}\) 到 \(\mathbf{Y}\) 中的）通过考虑映射 \(x \to \sum_{\lambda \in \mathbf{J}} u_\lambda(x)\)（对 L 的所有有限子集

J）以及这些映射在 \(\mathcal{F}_{\mathfrak{S}}(\mathbf{X};\mathbf{Y})\) 中关于 L 的有限子集的定向集的极限来定义（第 III 章，§ 5，第 1 号）。

\begin{enumerate}
\def\labelenumi{\arabic{enumi})}
\setcounter{enumi}{5}
\tightlist
\item
  由定义 1 与定义 2 立即得出：对每个 \(x \in \bigcup_{\mathbf{A} \in \mathfrak{S}} \mathbf{A}\)，映射 \(u \to v(x)\)（\(\mathcal{F}_{\mathfrak{S}}(\mathbf{X}; \mathbf{Y})\) 到 Y 中的）是一致连续的。因而特别地，若用 \(\overline{\mathbf{H}}\) 表示 \(\mathcal{F}_{\mathfrak{S}}(\mathbf{X}; \mathbf{Y})\) 的子集 H 的闭包，则有 \(\overline{\mathbf{H}}(x) \subset \overline{\mathbf{H}(x)}\)（对所有 \(x \in \bigcup_{\mathbf{A} \in \mathfrak{S}} \mathbf{A}\)）（第 I 章，§ 2，第 1 号，定理 1）。
\end{enumerate}

